\documentclass[10pt,a4paper]{amsart}
\usepackage{amsmath,amssymb,amsthm,graphicx}
\usepackage{verbatim}
\newenvironment{reference}{\comment}{\endcomment}
%\newenvironment{reference}{}{}
\newenvironment{slogan}{\comment}{\endcomment}
\newenvironment{history}{\comment}{\endcomment}

% For commutative diagrams we use Xy-pic
\usepackage[all]{xy}

% We use 2cell for 2-commutative diagrams.
\xyoption{2cell}
% Theorem environments.
%
\theoremstyle{plain}
\newtheorem{theorem}[subsection]{Theorem}
\newtheorem{proposition}[subsection]{Proposition}
\newtheorem{lemma}[subsection]{Lemma}

\theoremstyle{definition}
\newtheorem{definition}[subsection]{Definition}
\newtheorem{example}[subsection]{Example}
\newtheorem{exercise}[subsection]{Exercise}
\newtheorem{situation}[subsection]{Situation}

\theoremstyle{remark}
\newtheorem{remark}[subsection]{Remark}
\newtheorem{remarks}[subsection]{Remarks}

\numberwithin{equation}{subsection}

% Macros
%
\def\lim{\mathop{\mathrm{lim}}\nolimits}
\def\colim{\mathop{\mathrm{colim}}\nolimits}
\def\Spec{\mathop{\mathrm{Spec}}}
\def\Hom{\mathop{\mathrm{Hom}}\nolimits}
\def\Ext{\mathop{\mathrm{Ext}}\nolimits}
\def\SheafHom{\mathop{\mathcal{H}\!\mathit{om}}\nolimits}
\def\SheafExt{\mathop{\mathcal{E}\!\mathit{xt}}\nolimits}
\def\Sch{\mathit{Sch}}
\def\Mor{\mathop{\mathrm{Mor}}\nolimits}
\def\Ob{\mathop{\mathrm{Ob}}\nolimits}
\def\Sh{\mathop{\mathit{Sh}}\nolimits}
\def\NL{\mathop{N\!L}\nolimits}
\def\CH{\mathop{\mathrm{CH}}\nolimits}
\def\proetale{{pro\text{-}\acute{e}tale}}
\def\etale{{\acute{e}tale}}
\def\QCoh{\mathit{QCoh}}
\def\Ker{\mathop{\mathrm{Ker}}}
\def\Im{\mathop{\mathrm{Im}}}
\def\Coker{\mathop{\mathrm{Coker}}}
\def\Coim{\mathop{\mathrm{Coim}}}

% Boxtimes
%
\DeclareMathSymbol{\boxtimes}{\mathbin}{AMSa}{"02}

%
% Macros for moduli stacks/spaces
%
\def\QCohstack{\mathcal{QC}\!\mathit{oh}}
\def\Cohstack{\mathcal{C}\!\mathit{oh}}
\def\Spacesstack{\mathcal{S}\!\mathit{paces}}
\def\Quotfunctor{\mathrm{Quot}}
\def\Hilbfunctor{\mathrm{Hilb}}
\def\Curvesstack{\mathcal{C}\!\mathit{urves}}
\def\Polarizedstack{\mathcal{P}\!\mathit{olarized}}
\def\Complexesstack{\mathcal{C}\!\mathit{omplexes}}
% \Pic is the operator that assigns to X its picard group, usage \Pic(X)
% \Picardstack_{X/B} denotes the Picard stack of X over B
% \Picardfunctor_{X/B} denotes the Picard functor of X over B
\def\Pic{\mathop{\mathrm{Pic}}\nolimits}
\def\Picardstack{\mathcal{P}\!\mathit{ic}}
\def\Picardfunctor{\mathrm{Pic}}
\def\Deformationcategory{\mathcal{D}\!\mathit{ef}}
\usepackage{cleveref}
\crefname{originallemma}{original lemma}{original lemmas}
\Crefname{originallemma}{Original lemma}{Original lemmas}
\makeatletter
\newlabel{lemma-check-exactness-fibre-product}{{08N5}{source}{}{sourcepointer-1}{}}
\newlabel{lemma-check-exactness-fibre-product@cref}{{[originallemma][0][]08N5}{[1][0][]source}{}{}{}}
\makeatother
\title{Exactness of the induced sequences on kernels and cokernels in a commutative diagram}
\author{The Stacks Project Authors}
\date{stacks-project:homology.tex:local-snapshot-20261012; GNU FDL 1.2 or later}
\begin{document}
\maketitle
\noindent\textit{Editorial note.} All delivered mathematical text is the author's own, copied from the pinned Stacks Project chapter file homology.tex; no statement or proof step was rewritten, reordered or removed.\par\medskip
\noindent\textit{Editorial note.} The proof is the author's complete argument for the kernel statement, including the two applications of the exactness criterion; the closing clause 'and then (2) follows by duality' is the author's own routine duality remark on the part written out in full inside the counted span, and it is preserved verbatim.\par\medskip
\noindent\textit{Editorial note.} Every labelled reference occurring inside the counted spans is carried in assembly.external\_references as a non-load-bearing navigation pointer (same-chapter ones at their definition site, cross-chapter ones at the invocation inside the counted span), so no load-bearing mathematics is replaced by a stub. No bibliography citation occurs inside the counted spans.\par\medskip
\noindent\textit{Editorial note.} The same-name check for this card scanned every staged Stacks chapter of this repository -- all sources/stacks-*/pinned-original-source.tex files, the card's own chapter excluded -- and the labels found are recorded in dedup.same\_name\_cross\_chapter on the card.\par\medskip
\noindent\textit{Editorial note.} The preamble is the chapter's own native preamble excerpt already used for this source: the comment-type reference, slogan and history environments (native preamble lines 8-13), the Xy-pic setup (native preamble lines 14-18) and the native theorem declarations with the mathematical macros (native preamble lines 154-217, the actual last line of the file).\par\medskip
\noindent\textit{Allowed standard prerequisites (editorial).}\begin{itemize}
\item Testing exactness through a fibre product. Same chapter, Lemma lemma-check-exactness-fibre-product. Supplies the criterion that converts the required exactness statements into the existence of auxiliary objects and morphisms, applied twice inside the proof and once at its end.
\end{itemize}\medskip
\ifdefined\hypertarget\hypertarget{sourcepointer-1}{}\fi
\noindent\textit{Source reference.} Original lemma 08N5: sources/stacks-homology/pinned-original-source.tex, source lines 1152--1152. Navigation pointer invoked at proof line 1216; the counted span contains the pointer only and the cited result is not reproduced inside it, so the pointer is registered as a non-load-bearing navigation stub in the established form for Stacks material. No bibliography citation is involved.\par\medskip
\begin{lemma}
\label{lemma-exact-kernel-sequence}
Let $\mathcal{A}$ be an abelian category. Let
$$
\xymatrix{
x \ar[r]^f \ar[d]^\alpha &
y \ar[r]^g \ar[d]^\beta &
z \ar[d]^\gamma\\
u \ar[r]^k & v \ar[r]^l & w
}
$$
be a commutative diagram.
\begin{enumerate}
\item If the first row is exact and $k$ is a monomorphism, then the induced 
sequence $\Ker(\alpha) \to \Ker(\beta) \to \Ker(\gamma)$ 
is exact.
\item If the second row is exact and $g$ is an epimorphism, then the induced 
sequence
$\Coker(\alpha) \to \Coker(\beta) \to \Coker(\gamma)$ 
is exact.
\end{enumerate}
\end{lemma}

\begin{proof}
Suppose the first row is exact and $k$ is a monomorphism. Let 
$a:\Ker(\alpha)\to\Ker(\beta)$ and 
$b:\Ker(\beta)\to\Ker(\gamma)$ be the induced morphisms. 
Let $h:\Ker(\alpha)\to x$, $i:\Ker(\beta)\to y$ and 
$j:\Ker(\gamma)\to z$ be the canonical injections. As $j$ is 
a monomorphism we have $b\circ a=0$. Let $c:s\to\Ker(\beta)$ 
with $b\circ c=0$. Then, $g\circ i\circ c=j\circ b\circ c=0$. By 
Lemma \ref{lemma-check-exactness-fibre-product} there are an object $t$, an 
epimorphism $d:t\to s$ and a morphism $e:t\to x$ with 
$i\circ c\circ d=f\circ e$. Then, 
$k\circ \alpha\circ e=\beta\circ f\circ e=\beta\circ i\circ c\circ d=0$. 
As $k$ is a monomorphism we get $\alpha\circ e=0$. So, there exists 
$m:t\to\Ker(\alpha)$ with $h\circ m=e$. It follows 
$i\circ a\circ m=f\circ h\circ m=f\circ e=i\circ c\circ d$. 
As $i$ is a monomorphism we get $a\circ m=c\circ d$. Thus, 
Lemma \ref{lemma-check-exactness-fibre-product} implies (1), and then 
(2) follows by duality.
\end{proof}

\end{document}
